% --------------------------------------------------------------------------------------- %
%                   UOC - CSD - HY119 - Spring 2025 - Ass3 - csd5748 - :)                 %
% --------------------------------------------------------------------------------------- %


\documentclass[fleqn]{article}
\usepackage{amsmath, amssymb}
\setlength{\mathindent}{0pt}
\usepackage[utf8]{inputenc} 
\usepackage[english, greek]{babel}
\usepackage{subcaption}
\usepackage{hyperref}
\usepackage{xcolor}
\hypersetup
{
    hidelinks,  
    colorlinks=true,
    linkcolor=blue,
    urlcolor=blue 
}
\captionsetup[subfigure]{labelformat=empty}
\usepackage{amsmath}
\usepackage{amsfonts}
\usepackage{amssymb}
\usepackage{array}
\usepackage{graphicx}
\usepackage[a4paper, left=1.1cm, right=1.1cm, top=1.3cm, bottom=2cm]{geometry}
\usepackage[absolute, overlay]{textpos}
\usepackage{placeins}
\usepackage{listings}
\usepackage{xcolor}
\setlength{\jot}{0pt}
\lstdefinestyle{mypython}
{
    language=Python,
    basicstyle=\ttfamily\small,
    keywordstyle=\color{blue},
    stringstyle=\color{red},
    commentstyle=\color{green!50!black},
    numbers=left,
    numberstyle=\tiny\color{gray},
    stepnumber=1,
    frame=single,
    backgroundcolor=\color{gray!10},
    breaklines=true
}
\usepackage{cancel}
\setlength{\mathsurround}{0.1cm}
\setlength{\jot}{0.1cm}
\newcommand{\skato}{\hspace{0.1cm}}
\newcommand{\skata}{\hspace{0.321cm}}
\begin{document}
\selectlanguage{greek}



% ---------------------------------------------------------------------------------------



\begin{center}
\boxed{\textlatin{csd5748}}
\end{center}
\vspace{-0.5cm}
\begin{center}
\setlength{\fboxsep}{20pt}
\noindent{\boxed{{{\textbf{\Large{ΗΥ-119 - Γραμμική Άλγεβρα -  Ασκήσεις 3.}}}}}}
\end{center}
\vspace{0cm}


% ---------------------------------------------------------------------------------------



\vspace{1.23456789cm}
\begin{center}
    {{\textbf{{[ 1\textsuperscript{η} Άσκηση ]}}}}
\end{center}
\vspace{0.01cm}

\[
\textbf{Έχουμε:} \hspace{0.54321cm} 
A \skato = \skato
\left(
\begin{matrix}
& 1 & 1 & 3 &  \\
& 0 & 4 & 6 &  \\
& 1 & 5 & 8 &
\end{matrix} 
\right)
\]

\[
\textbf{\textlatin{i)} Ξέρουμε:} \hspace{0.54321cm} 
H \skato = \skato
\left(
\begin{matrix}
& a & b & c &  \\
& d & e & f &  \\
& g & h & i &
\end{matrix} 
\right)
\implies
\det(H) = (aei+bfg+cdh)-(ceg+afh+bdi)
\]

\[
\textbf{Έτσι:} \hspace{0.54321cm}  
\det(A) = ((1)(4)(8)+(1)(6)(1)+(3)(0)(5))-((3)(4)(1)+(1)(6)(5)+(1)(0)(8)) = 38 - 42 = -4
\]



\[
\textbf{\textlatin{ii)} Ξέρουμε:} \hspace{0.54321cm} 
H_{3 \times 3}
\overset{Gauss}{\rightsquigarrow}
\begin{bmatrix}
    &d_1&*&*& \\ &0&d_2&*& \\ &0&0&d_3&
\end{bmatrix}
\implies
\det(H) = d_1 \cdot d_2 \cdot d_3
\]

\[
\textbf{Έτσι:} \hspace{0.54321cm}  
A = \begin{bmatrix}
& 1 & 1 & 3 &  \\
& 0 & 4 & 6 &  \\
& 1 & 5 & 8 &
\end{bmatrix}
\overset{r_3 \leftarrow r_3 - r_1}{\rightsquigarrow}
\begin{bmatrix}
& 1 & 1 & 3 &  \\
& 0 & 4 & 6 &  \\
& 0 & 4 & 5 &
\end{bmatrix}
\overset{r_3 \leftarrow r_3 - r_2}{\rightsquigarrow}
\begin{bmatrix}
& 1 & 1 & 3 &  \\
& 0 & 4 & 6 &  \\
& 0 & 0 & -1 &
\end{bmatrix}
\implies
\det(H) = (1)(4)(-1) = -4
\]

\[
\textbf{Τελικά:} 
\hspace{0.54321cm} 
\boxed{ \
\det(A) = -4 \
}
\]


% -----


\vspace{1.23456789cm}
\begin{center}
    {{\textbf{{[ 2\textsuperscript{η} Άσκηση ]}}}}
\end{center}
\vspace{0.01cm}



\[
\textbf{Ξέρουμε:} \hspace{0.54321cm} 
\det(A) \hspace{0.321cm} =  
\sum_{P_i(i_1 \dots i_n)} \ \sigma(P_i) \cdot a_{1 i_1} \cdot a_{2 i_2} \cdot \ \cdots \ \cdot a_{n i_n} 
\]

\[
\textbf{Όμως:} \hspace{0.54321cm} 
( A^T )_{ij} \hspace{0.321cm} = \hspace{0.321cm} A_{ji}
\]

\[
\textbf{Έτσι:} \hspace{0.54321cm} 
\det(M) \hspace{0.321cm} =  
\sum_{P_i(i_1 \dots i_n)} \sigma(P_i) \cdot m^T_{1 i_1} \cdot m^T_{2 i_2} \cdot \cdots \cdot m^T_{n i_n}
\hspace{0.321cm}
=  
\sum_{P_i(i_1 \dots i_n)} \sigma(P_i) \cdot m_{i_11} \cdot m_{i_22} \cdot \cdots \cdot m_{i_nn}
\hspace{0.321cm}
=
\hspace{0.321cm}
\det(M^T)
\]

\[
\textbf{Τελικά:} 
\hspace{0.54321cm} 
\boxed{
\ \det(M) = \det(M^T) \
}
\]


% ------


\vspace{3.456789cm}
\begin{center}
    {{\textbf{{[ 3\textsuperscript{η} Άσκηση ]}}}}
\end{center}
\vspace{-0.1cm}



\[
\textbf{Έχουμε:} \quad
[1, 5, 0]^T, \quad [0, 2, 1]^T, \quad [3, 1, 8]^T
\quad \implies \quad
\vec{u}_1 =
\begin{bmatrix}
1 \\ 5 \\ 0
\end{bmatrix}, \quad
\vec{u}_2 =
\begin{bmatrix}
0 \\ 2 \\ 1
\end{bmatrix}, \quad
\vec{u}_3 =
\begin{bmatrix}
3 \\ 1 \\ 8
\end{bmatrix}
\quad \implies \quad
A = 
\begin{bmatrix}
1 & 0 & 3 \\
5 & 2 & 1 \\
0 & 1 & 8
\end{bmatrix}
\]

\[
\textbf{Ορίζουσα:} \quad
\det(A) =
1 \cdot 
\begin{vmatrix}
2 & 1 \\
1 & 8
\end{vmatrix}
- 
0 \cdot 
\begin{vmatrix}
5 & 1 \\
0 & 8
\end{vmatrix}
+ 
3 \cdot 
\begin{vmatrix}
5 & 2 \\
0 & 1
\end{vmatrix}
=
1 \cdot (2 \cdot 8 - 1 \cdot 1) - 0 + 3 \cdot (5 \cdot 1 - 0 \cdot 2)
\ = \ 30 \ \ne \ 0
\]   

\[
\textbf{Τελικά, τα διανύσματα} \ \boxed{\textbf{είναι γραμμικά ανεξάρτητα.}}
\]


 
% -----



\vspace{1cm}
\begin{center}
    {{\textbf{{[ 4\textsuperscript{η} Άσκηση ]}}}}
\end{center} 
\vspace{-0.25cm}

\[
\textbf{Έχουμε:} \hspace{0.54321cm} 
B = 
\begin{bmatrix}
1 & 2 & 3 \\
0 & 4 & 0 \\
0 & 0 & 5
\end{bmatrix}
\]

\[
\textbf{Ξέρουμε:} \hspace{0.54321cm}
\text{Ανάπτηγμα σε συμπαράγοντες:} \quad C_{ij} = \sigma_{ij} \cdot \det(A^{ij}) = (-1)^{i+j} \cdot \det(A - A_{ij})
\]

\[
\textbf{Είναι:} \ \ 
\det(B) = \beta_{11}C_{11} + \beta_{21}C_{21} + \beta_{31}C_{31} = (1)(-1)^2(4\cdot5-0\cdot0) + (0)(-1)^3(2\cdot5-3\cdot0) + (0)(-1)^4(2\cdot0-3\cdot4) = 20 \neq 0 \Rightarrow \exists \ B^{-1}
\]

\[
\textbf{Ξέρουμε:} \hspace{0.54321cm}
\text{Αντίστροφος πίνακας:} \quad B^{-1} \ = \ \frac{1}{\det(B)} \ \cdot \ C^T
\]

\[
\textbf{Είναι:} \quad \
C \ = \ 
\begin{bmatrix}
    & (+1) \cdot \begin{vmatrix} 4&0\\0&5  \end{vmatrix}
    & \ \ (-1) \cdot \begin{vmatrix} 0&0\\0&5  \end{vmatrix} \ \
    & (+1) \cdot \begin{vmatrix} 0&4\\0&0  \end{vmatrix}
    & \\
    & \\
    & (-1) \cdot \begin{vmatrix} 2&3\\0&5  \end{vmatrix}
    & \ \ (+1) \cdot \begin{vmatrix} 1&3\\0&5  \end{vmatrix} \ \
    & (-1) \cdot \begin{vmatrix} 1&2\\0&0  \end{vmatrix}
    & \\
    & \\
    & (+1) \cdot \begin{vmatrix} 2&3\\4&0  \end{vmatrix}
    & \ \ (-1) \cdot \begin{vmatrix} 1&3\\0&0  \end{vmatrix} \ \
    & (+1) \cdot \begin{vmatrix} 1&2\\0&4  \end{vmatrix}
\end{bmatrix}
\ = \
\begin{bmatrix}
     20 & 0 & 0  \\
     -10 & 5 & 0  \\
     -12 & 0 & 4 
\end{bmatrix}
\ \ \Rightarrow \ \ \
C_{3 \times 3}^T \ = \
\begin{bmatrix}
    & 20 & -10 & -12 & \\
    & 0 & 5 & 0 & \\
    & 0 & 0 & 4 &
\end{bmatrix}
\]

\[
\textbf{Τελικά:} \quad \ \ 
B^{-1} \ = \
\frac{1}{20} \ \cdot \ 
\begin{bmatrix}
    & 20 & -10 & -12 & \\
    & 0 & 5 & 0 & \\
    & 0 & 0 & 4 &
\end{bmatrix}
\ = \
\boxed{ 
\begin{bmatrix}
    & 1 & -1/2 & -3/5 & \\
    & 0 & 1/4 & 0 & \\
    & 0 & 0 & 1/5 &
\end{bmatrix}
\ = \ 
B^{-1}_{3 \times 3}
\ }
\]


% -----


\vspace{1cm}
\begin{center}
    {{\textbf{{[ 5\textsuperscript{η} Άσκηση ]}}}}
\end{center}
\vspace{0.01cm}

\[
\textbf{Έχουμε:} \hspace{0.54321cm} 
\lambda_1 = 1, \ \ \lambda_2 = 2, \ \ \lambda_3 = 4 \ \ \text{ιδιοτιμές του πίνακα.} \ 
\overset{\dots \ }{\implies} \ C \ = \ 3 \times 3.
\]

\[
\textbf{Ξέρουμε:} \quad \ 
\text{Ίχνος πίνακα:} \quad 
tr(C) \ = \ \sum_{i=1}^n ( \lambda_i) \ \overset{*}{\implies} \ tr(C^2) \ = \ \sum_{i=1}^n ( \lambda_i^2)
\ \overset{\dots \ }{\implies} \ tr(C^2) \ = \ (\lambda_1)^2 \ + \ (\lambda_2)^2 \ + \ (\lambda_3)^2 \ = \ \boxed{ \ 21 \ }_\alpha
\]

\[
\textbf{Επιπλέον:} \quad \ 
\det((C^{-1})^T) \ \overset{\text{Ιδ.}}{=} \ \det(C^{-1}) \ \overset{\text{Ιδ.}}{=} \ \frac{1}{\det(C)} \ = \ \frac{1}{\prod_{i=1}^{n} (\lambda_i)}
\ \overset{\dots \ }{\implies} \ 
\det((C^{-1})^T) \ = \ \frac{1}{(1)(2)(4)} \ = \ \boxed{ \ \frac{1}{8} \ }_\beta
\]

 
 

 

% ---------------------------------------------------------------------------------------

\newpage

\begin{center}
    \vspace*{\fill}
    \setlength{\fboxsep}{10pt}
    \boxed{\textbf{\LARGE{ΤΕΛΟΣ}}}
    \vspace*{\fill}
\end{center}

\end{document} 

% ---------------------------------------------------------------------------------------



 
% The end ...